\chapter{Introduction}
\label{chap:final-introduction}
This chapter introduces the clinical and methodological problem addressed in this dissertation: metal artifact reduction in computed tomography. It first situates the development of X-ray imaging and\acrshort{ct}, then describes the clinical and technical consequences of metal artifacts. The chapter closes by motivating a reproducible, safety-oriented evaluation of image-domain deep learning methods; the study objectives and hypotheses are specified in the following chapter.

In 1895, while repeating an experiment on cathode rays with a modified Crookes tube, the German physicist Wilhelm Conrad Röntgen observed a new, highly penetrating form of radiation and used it to produce the first radiograph of the human body. Later naming the phenomenon "X-rays" \cite{November81895,okunoFisicaRadiacoes2010}. The discovery spread rapidly: within months, X-rays were already being used for medical diagnosis, such as in a fracture radiograph produced in the United States in early 1896 \cite{November81895}. Since then, the technique has undergone successive modernizations, diversifying into multiple modalities and devices and consolidating itself as an indispensable tool in modern healthcare \cite{howellEARLYCLINICALUSE2016, taylorIlluminatingHealthEvolution2023}.

From these successive modernizations emerged, in 1971, the first computerized tomography scanner. This technology is based on the principle of x-rays attenuation, in which the amount of radiation absorbed or scattered by the material is related to its density \cite{valiUnveilingMarvelsComputed2024, noorOverviewComputedTomography2023}.

\acrfull{ct} was the first medical imaging modality to provide quantitative information about tissue attenuation, something that was not possible with conventional radiography \cite{schulzHowCTHappened2021}. Furthermore, allowing sectional visualization of the human body, it eliminated the overlap of anatomical structures inherent to two-dimensional images, enabling precise localization and assessment of the extent of lesions \cite{schulzHowCTHappened2021}. 

However, these advances brought a significant increase in technological complexity. The acquisition of multiple projections, combined with computational reconstruction and digital image processing, introduced new forms of degradation that were nonexistent or of little relevance in conventional radiography. Thus, computed tomography became particularly susceptible to the presence of artifacts, which can compromise the diagnostic quality and clinical reliability of images \cite{schulzHowCTHappened2021}.
 
Various types of artifacts can be observed in computed tomography, including noise, beam hardening, motion, ring artifacts, and metal artifacts, which can obscure or simulate pathologies \cite{boasCTArtifactsCauses2012}. The causes of these artifacts are varied, ranging from detector failures or miscalibration, statistical limitations associated with low photon counts, to physical phenomena such as beam hardening and scattering, which produce dark streaks between high-attenuation objects such as metal, bone, or iodinated contrast media \cite{schulzHowCTHappened2021}.

Specifically, streak artifacts caused by metallic implants are observed in approximately 21\% of computed tomography scans, being considered extremely common in clinical practice \cite{schulzHowCTHappened2021}. The high attenuation or even total blockage of X-rays by these implants prevents the radiation from adequately reaching the detectors, resulting in low-quality projection data and compromising image reconstruction \cite{kleberAdvancementsSupervisedDeep2024}.

The increased use of metallic implants, driven by the evolution of healthcare, has intensified concerns regarding the negative impact of these artifacts on image quality and the reliability of diagnostic analyses \cite{kleberAdvancementsSupervisedDeep2024}. Traditionally, various\acrfull{mar} algorithms have been proposed, including image inpainting techniques, sinogram inpainting, and\acrfull{mbir} or combinations of these approaches \cite{puvanasunthararajahApplicationMetalArtifact2021}.

Although analytical methods such as\acrfull{fbp} have dominated tomographic reconstruction for decades due to their computational efficiency and conceptual simplicity, their limitations become evident under non-ideal clinical conditions, requiring higher radiation doses to ensure image quality \cite{sahuTransformingCTImaging2025}.\acrshort{mbir} methods emerged as an alternative to overcome these limitations, at the cost of higher computational demand and the introduction of artificial textures in the reconstructed image \cite{sahuTransformingCTImaging2025}.

Recently, deep learning-based approaches have been widely investigated for artifact correction in computed tomography, demonstrating superior performance and greater clinical applicability compared to conventional methods \cite{alkhodariEnhancingCCTAImage2025}. In particular,\acrfullpl{cnn} have been proposed for metal artifact reduction, enabling noise suppression and preservation of anatomical and soft tissue details \cite{alkhodariEnhancingCCTAImage2025}.


\section{Clinical and Technical Context}

A diagnostic mistake can have serious consequences for the patient, including treatment delays, inadequate or unnecessary treatments, and even life-threatening risks \cite{alkhodariEnhancingCCTAImage2025}. Artifacts can degrade an image to the point of making it unusable for diagnosis, in which case the scan may need to be repeated, exposing the patient to further radiation. When the artifact instead obscures nearby soft tissue rather than making the whole image unusable, it can still lead to an incorrect diagnosis or, in radiotherapy planning, an incorrect radiation dose calculation \cite{alzainCommonComputedTomography2021}.

Therefore, the reduction of artifacts in\acrshort{ct} is a research area with significant clinical importance, and the development of effective methods to reduce artifacts gains high value in medical practice.

\section{Work Motivation}

Although studies using deep learning for metal artifact reduction are not new, a lot of the work published to date does not present a rigorous and detailed evaluation of the performance of the proposed models. In particular, the evaluation of the performance of models on external datasets (which were not used during model training) or evaluation of images by qualified professionals is often lacking. Moreover, the training of supervised models depends on reference images, which are difficult to obtain in a clinical context, creating a dependency on synthetic or simulated images. This gap becomes even more relevant when considering that the performance of models trained on synthetic images may not be representative of their performance on real images. This mismatch can lead to overestimation of their effectiveness and limit their clinical applicability.

In this sense, the present study proposes a more rigorous and detailed evaluation of the performance of deep learning models for this purpose.




